
%%============================================================
\section{Methods}
%%============================================================

\subsection{Study identification and matched-set extraction}

Ten seed papers comparing chelators on a common vector were expanded through their Semantic Scholar citation and reference networks, yielding 681 unique DOIs. Titles and abstracts were screened for two or more distinct chelators from a controlled vocabulary of 28 chelator cores, a measured affinity endpoint, a targeting vector, and comparison language, giving 42 candidate studies of which 38 (90\%) were open access. Full texts were retrieved as structured XML through the Europe PMC REST API.

Affinity tables were read and transcribed by hand; no value was obtained by automated parsing, interpolation or estimation. Each record carries its source DOI, vector, chelator, metal, endpoint, value, standard deviation and a provenance note. Values that a source reproduced from earlier work were identified by their footnotes and excluded as duplicates---three pentixafor studies\cite{Poschenrieder2016a,Poschenrieder2016b,Poschenrieder2017} share one underlying dataset.

The analysis unit is the \emph{matched set}: one paper, one vector, one metal, one endpoint, with the chelator varying, so that assay, cell line, radioligand and vector are constant. The final dataset comprises 106 values from ten studies,\cite{Poschenrieder2016a,Poschenrieder2017,Mitran2019,Fischer2025,Weineisen2014,Sinnes2020,Basaco2024a,Basaco2024b,Rinne2019,Rinne2020} of which 95 fall in 32 matched sets spanning five target families (CXCR4, GRPR, PSMA, SSTR, affibody/DARPin), 13 vectors, nine chelators and ten metal states. Endpoints were IC$_{50}$ (71), $K_\text{D}$ (18) and $K_\text{i}$ (6); set sizes ranged from two to eight chelator arms. Values were log$_{10}$-transformed and centred within their matched set.

\subsection{Descriptors}

\emph{Classical (10).} Net formal charge of the metal--chelate unit in the conjugated form (one arm consumed by the amide bond to the vector), number of free carboxylates, macrocycle nitrogen count and ring size, denticity, macrocyclic flag, metal formal charge, Shannon ionic radius, a metal-present flag, and a cavity-fit mismatch term.

\emph{Quantum-chemical (12).} For every (chelator, metal) combination in the dataset the conjugated ligand was built from SMILES with one arm capped as an $N$-methylamide, the metal threaded into the cavity using explicit dative bonds before 3-D embedding, and the complex relaxed with GFN2-xTB\cite{Bannwarth2019,Bannwarth2020} at the correct total charge and spin with ALPB implicit solvation in water. These are charged species---the conjugated complexes carry formal charges from $-4$ to $+1$---and a gas-phase treatment of them is not defensible. Five starting structures were used per complex (four fresh conformer seeds plus a re-relaxation of a gas-phase optimum) and the lowest-energy structure with a chemically valid coordination mode kept.

Complexation energies referenced to the bare metal ion are dominated by Coulombic neutralisation rather than chemistry: computed that way they span $-691$ to $-2257$\,kcal\,mol$^{-1}$ and correlate with the metal's formal charge at $r = -0.914$, and implicit solvation alone does not repair this ($r = -0.884$). We therefore used a \emph{ligand-exchange energy} referenced to DOTA at the same metal, $\Delta\Delta E(\mathrm{L};\mathrm{M}) = [E(\mathrm{ML}) - E(\mathrm{L})] - [E(\mathrm{M}\cdot\mathrm{DOTA}) - E(\mathrm{DOTA})]$, in which the ion term cancels exactly. It spans $-35.7$ to $+20.6$\,kcal\,mol$^{-1}$ and its correlation with metal charge falls to $r = -0.199$ ($p = 0.320$).

Coordination geometry was verified against the known coordination mode of each chelator rather than assumed, giving 27/27 verified complexes (Ga--DOTA N$_4$O$_3$ with Ga--O 1.84--1.94\,\AA{} and Ga--N 2.26--2.54\,\AA; Ga--NODAGA N$_3$O$_3$; Lu--DOTA N$_4$O$_4$). Solvation changed the structures materially---four of 27 coordination numbers shifted, mean M--N contracted by 0.10\,\AA, and the mean dipole moment rose from 16.0 to 25.9\,D.

Every log$D_{7.4}$ or log$P$ value reported alongside the affinity data in the same tables (38 values from four of the ten studies) was transcribed and tested as a confounder.

\subsection{Models and validation}

All models predicted the within-set deviation. Comparators were: predicting zero; a one-hot chelator encoding, which is the model implied by a universal chelator ranking; classical descriptors; and quantum-chemical descriptors, alone and combined. Continuous features were standardised and fitted with ridge regression.

Validation was leave-one-target-family-out. Random splits are invalid here: matched sets contribute several rows sharing a vector, so a random split places one arm of a construct in training and another in test, leaking the set mean. Transfer was additionally evaluated with held-out chelator, vector, paper and metal. Confidence intervals come from 10\,000 bootstrap resamples and ten random seeds were averaged. A within-set label-permutation test (2000 permutations) provides the null distribution of LOPO $R^2$ under no transferable information. Normality was assessed by Shapiro--Wilk ($W = 0.982$, $p = 0.233$); paired contrasts used Wilcoxon signed-rank with Cohen's $d_z$.

Three pre-specified sanity checks were applied before interpretation---monotonicity of the response against complex charge, correctness of the within-set centring, and consistency between the IC$_{50}$-only subset and the full set. All three passed ($\rho = -0.208$, $p = 0.043$ in the hypothesised direction; maximum absolute within-set mean $1.7\times10^{-16}$; IC$_{50}$-only $\rho = -0.243$).

%%============================================================
\section*{Conflicts of interest}
%%============================================================
There are no conflicts to declare.

%%============================================================
\section*{Data availability}
%%============================================================
The data supporting this article are available in full. The matched affinity dataset with per-row DOI provenance, the derived analysis table, the quantum-chemical descriptors and optimised Cartesian geometries of all 32 metal--chelate complexes, the screened candidate pool and the complete analysis pipeline are deposited at [repository DOI to be inserted at acceptance] under CC-BY. Every affinity value was transcribed by hand from an open-access full text and carries its source DOI; no value was parsed automatically, interpolated, estimated or generated. Running the pipeline from the deposited files reproduces every number and figure in this article.

%%============================================================
\section*{Author contributions}
%%============================================================
J.K. and J.Y. contributed equally. J.K.: quantum-chemical calculations, methodology, validation, writing---review and editing. J.Y.: data curation, investigation, validation, writing---review and editing. K.K.: conceptualization, formal analysis, methodology, software, visualization, writing---original draft, supervision. D.L.: conceptualization, interpretation, writing---review and editing, supervision. All authors approved the final manuscript. [Per-author allocation to be confirmed by the co-authors before submission.]

%%============================================================
\section*{Acknowledgements}
%%============================================================
[To be completed.]

%%============================================================
\balance
\renewcommand\refname{Notes and references}
\bibliographystyle{rsc}
\bibliography{references}

\end{document}
